\subsection{Classifying space for a discrete group}

Let G be a topological group; denote its identity element by e. Let G* be the set of maps \(h\colon[0,1[ \to G\) for which there exists a finite sequence \((t_{i})_{0\leqslant i\leqslant n}\) with \(0=t_{0}<t_{1}<\cdots<t_{n}=1\) such that h is constant on the intervals \([t_{i-1},t_{i}[\) for \(1\leqslant i\leqslant n\). Such a sequence will be said to be adapted to h. For every finite subset of G*, there exists a sequence adapted to each of its elements.

The set \(G^{*}\) is a subgroup of \(G^{[0,1[}\). We denote by \(e^{*}\) its identity element; the inverse of an element \(h \in G^{*}\) is the map \(t \mapsto h(t)^{-1}\), denoted \(h^{-1}\).

Let V be a neighborhood of \(e\) in G. Let \(h \in \mathbf{G}^*\) and let \((t_i)_{0 \leqslant i \leqslant n}\) be a sequence adapted to \(h\). The set of elements \(t \in [0,1[\) such that \(h(t) \notin \mathbf{V}\) is the union of some of the intervals \([t_{i-1}, t_i[\); the sum of the lengths \(t_i - t_{i-1}\) of these intervals does not depend on the sequence \((t_i)_{0 \leqslant i \leqslant n}\) chosen; denote it by \(p_{\mathrm{V}}(h)\). For every real number \(\varepsilon\) such that \(\varepsilon > 0\), denote by \(\mathrm{V}_{\varepsilon}^*\) the set of \(h \in \mathbf{G}^*\) such that \(p_{\mathrm{V}}(h) < \varepsilon\).

PROPOSITION 5. --- There exists a unique topology on G* compatible with its group structure for which the sets V* form a base of neighborhoods of the identity element.

Let us verify that the sets \(V_{\varepsilon}^{*}\) satisfy the axioms \((\mathrm{GV}_{\mathrm{I}}^{\prime})\) , \((\mathrm{GV}_{\mathrm{II}}^{\prime})\) and \((\mathrm{GV}_{\mathrm{III}}^{\prime})\) from TG, III, p. 4.

Let V be a neighborhood of \(e\) in G and let \(\varepsilon\) be a strictly positive real number. Let W be a neighborhood of \(e\) in G such that \(\mathrm{W} \cdot \mathrm{W} \subset \mathrm{V}\). Let \(h\) and \(h'\) be elements of \(\mathbf{G}^*\). If \(t \in [0,1[\) is such that \(h(t)h'(t) \notin \mathrm{V}\), then \(h(t) \notin \mathrm{W}\) or \(h'(t) \notin \mathrm{W}\). It follows that \(p_{\mathrm{V}}(hh') \leqslant p_{\mathrm{W}}(h) + p_{\mathrm{W}}(h')\). Consequently, \(\mathrm{W}_{\varepsilon/2}^* \cdot \mathrm{W}_{\varepsilon/2}^* \subset \mathrm{V}_{\varepsilon}^*\), which proves the axiom ( \(\mathrm{GV}_I'\) ).

Let W be a neighborhood of e in G such that \(W^{-1} \subset V\) . Then \((\mathrm{W}_{\varepsilon}^{*})^{-1} = (\mathrm{W}^{-1})_{\varepsilon}^{*} \subset \mathrm{V}_{\varepsilon}^{*}\) , whence the axiom \((\mathrm{GV}_{\mathrm{II}}^{\prime})\) .

Let \(k\) be an element of \(\mathbf{G}^*\); since the function \(k\) takes only a finite number of values, there exists a neighborhood W of \(e\) in G such that \(k(t)\mathrm{W}k(t)^{-1}\subset \mathrm{V}\) for all \(t\in [0,1[\). Let \(h\) be an element of \(\mathbf{G}^*\). For \(t\in [0,1[\), if \(k(t)h(t)k(t)^{-1}\not\in\mathrm{V}\), then \(h(t)\not\in\mathrm{W}\). Consequently, \(p_{\mathrm{V}}(khk^{-1})\leqslant p_{\mathrm{W}}(h)\). This proves that \(k\mathrm{W}_{\varepsilon}^{*}k^{-1}\subset\mathrm{V}_{\varepsilon}^{*}\), whence the axiom ( \(\mathrm{GV}_{\mathrm{III}}'\) ).

PROPOSITION 6. --- The space G* is contractible and locally contractible at each of its points.

For \(h \in \mathbf{G}^*\) and \(t \in \mathbf{I}\), denote by \(\sigma(h, t)\) the map from [0,1[ into G given by \(\sigma(h, t)(x) = h(x)\) if \(0 \leqslant x < t\) and \(\sigma(h, t) = e\) otherwise.

Let us show that the map \(\sigma\colon\mathbf{G}^{*}\times\mathbf{I}\to\mathbf{G}^{*}\) thus defined is continuous. Indeed, let \(k\in\mathbf{G}^{*}, u\in\mathbf{I}\), let V be a neighborhood of \(e\) in G and let \(\varepsilon\) be a strictly positive real number. The element \(\sigma(h,t)\sigma(k,u)^{-1}\) of \(\mathbf{G}^{*}\) is the map \(f\) from [0,1[ into G given by

\[
f (x) = \left\{ \begin{array}{l l} h (x) k (x) ^ {- 1} & \text { if } 0 \leqslant x <   \min (t, u); \\ h (x) & \text { if } u \leqslant x <   t; \\ k (x) ^ {- 1} & \text { if } t \leqslant x <   u; \\ e & \text { otherwise. } \end{array} \right.
\]

Therefore,

\[
p _ {\mathrm{V}} (\sigma (h, t) \sigma (k, u) ^ {- 1}) \leqslant p _ {\mathrm{V}} (h k ^ {- 1}) + | t - u |;
\]

in other words, for \(\sigma(h,t) \in V_{\varepsilon}^{*}\sigma(k,u)\), it suffices that one have \(|t-u| \leqslant \frac{\varepsilon}{2}\) and \(h \in V_{\varepsilon/2}^{*}k\), which proves the continuity of \(\sigma\) at \((k,u)\). For all \(h \in G^{*}\), \(\sigma(h,0)\) is the constant map with image \(\{e\}\), while \(\sigma(h,1)=h\). Moreover, \(\sigma(e,t)=e\) for all \(t \in I\). Consequently, \(\sigma\) is a pointed homotopy at \(e \in G^{*}\) relating the constant map with image \(\{e\}\) to the identity map of \(G^{*}\). This proves that \(G^{*}\) is contractible to \(e^{*}\).

Furthermore, for every neighborhood V of e in G and every real number \(\varepsilon > 0\), we have \(\sigma(\mathrm{V}_{\varepsilon}^{*} \times \mathbf{I}) \subset \mathrm{V}_{\varepsilon}^{*}\). Consequently, \(V_{\varepsilon}^{*}\) is also contractible to \(e^{*} \in G^{*}\), so that \(G^{*}\) is locally contractible at \(e^{*}\).

Since G* is a topological group, it is contractible and locally contractible at each of its points.

Let \(\iota\) be the map from G to \(\mathrm{G}^*\) which associates with \(g\in \mathrm{G}\) the constant map with image \(\{g\}\) from [0,1[ into G. The map \(\iota\) is an injective group homomorphism. Let V be a neighborhood of \(e\) in G and let \(\varepsilon\) be a strictly positive real number. We have \(\iota^{-1}(\mathrm{V}_{\varepsilon}^{*}) = \mathrm{V}\) if \(\varepsilon \leqslant 1\) and \(\iota^{-1}(\mathrm{V}_{\varepsilon}^{*}) = \mathrm{G}\) otherwise. The inverse image of a neighborhood of the identity element of \(\mathrm{G}^*\) is a neighborhood of the identity element of G, whence the continuity of \(\iota\). Moreover, \(\iota(\mathrm{V}) = \mathrm{V}_1^* \cap \iota(\mathrm{G})\) for every neighborhood V of \(e\) in G. Consequently, \(\iota\) defines an isomorphism of topological groups from G onto its image.

Remark 1. --- If G is a separated topological group, \(\iota(\mathrm{G})\) is closed in \(\mathrm{G}^*\) .

Indeed, let \(h \in \mathbf{G}^*\) such that \(h \notin \iota(\mathbf{G})\), let \((t_i)_{0 \leqslant i \leqslant n}\) be a sequence adapted to \(h\), and set \(\varepsilon = \min_{1 \leqslant i \leqslant n} (t_i - t_{i-1})\). Let \(V\) be a neighborhood of \(e\) in \(G\) such that \(h(t_i)^{-1}h(t_j) \notin V\), for every pair \((i,j)\) of integers such that \(0 \leqslant i, j \leqslant n-1\) and \(h(t_i) \neq h(t_j)\); such a neighborhood exists because \(G\) is separated. Let \(W\) be a neighborhood of \(e\) in \(G\) such that \(W \cdot W^{-1} \subset V\). Let us prove that then \(hW_\varepsilon^*\) does not meet \(\iota(G)\).

Let us reason by contradiction. Let \(f\) be an element of \(W_{\varepsilon}^{*}\) and \(g\) be an element of \(G\) such that \(hf = \iota(g)\). We therefore have \(f(t) = h(t)^{-1}g\) for all \(t \in [0,1[\), so that the sequence \((t_i)\) is also adapted to the function \(f\). If the value taken by \(f\) on the interval \([t_{i-1}, t_i[\) does not belong to \(W\), then \(t_i - t_{i-1} < \varepsilon\) since \(f \in W_{\varepsilon}^{*}\). But this inequality is impossible by the definition of \(\varepsilon\); hence \(f(t) \in W\) for all \(t \in [0,1[\). Let then \(i\) and \(j\) be elements of \(\{0, \ldots, n-1\}\) such that \(h(t_i) \neq h(t_j)\); we have

\[
h (t _ {i}) ^ {- 1} h (t _ {j}) = f (t _ {i}) g ^ {- 1} g f (t _ {j}) ^ {- 1} = f (t _ {i}) f (t _ {j}) ^ {- 1} \in \mathrm{W} \cdot \mathrm{W} ^ {- 1},
\]

which contradicts the choice of W. The subgroup \(\iota(G)\) of \(G^{*}\) is therefore closed.

Suppose that G is a metrizable topological group and let d be a metric on G that defines its topology. Let h and \(h' \in G^{*}\) and let \((t_{i})_{0 \leqslant i \leqslant n}\) be a sequence adapted to h and \(h'\). The real number

\[
\sum_ {i = 1} ^ {n} (t _ {i} - t _ {i - 1}) d \bigl (h (t _ {i - 1}), h ^ {\prime} (t _ {i - 1}) \bigr)
\]

does not depend on the sequence \((t_{i})\) chosen; denote it by \(d^{*}(h,h')\) .

Let us prove that \(d^{*}\) is a distance on \(\mathbf{G}^{*}\). We have \(d^{*}(h,h^{\prime}) = d^{*}(h^{\prime},h)\) for \(h, h^{\prime} \in \mathbf{G}^{*}\), and \(d^{*}(h,h) = 0\) for all \(h \in \mathbf{G}^{*}\). Conversely, let \(h\) and \(h^{\prime}\) be elements of \(\mathbf{G}^{*}\) such that \(d^{*}(h,h^{\prime}) = 0\). Let \((t_{i})_{0\leqslant i\leqslant n}\) be a sequence adapted to \(h\) and \(h'\). Since

\[
0 = d ^ {*} (h, h ^ {\prime}) = \sum_ {i = 1} ^ {n} \left(t _ {i} - t _ {i - 1}\right) d \left(h \left(t _ {i - 1}\right), h ^ {\prime} \left(t _ {i - 1}\right)\right)
\]

and since all terms of this sum are positive or zero, we have \(d(h(t_{i-1}), h'(t_{i-1})) = 0\) for all \(i \in \{1, \ldots, n\}\) , whence \(h = h'\) . Finally, let \(h\) , \(h'\) , \(h''\) be elements of \(\mathbf{G}^*\) and let \((t_i)_{0 \leqslant i \leqslant n}\) a sequence adapted to each of them. Then,

\[
\begin{array}{l} d ^ {*} (h, h ^ {\prime \prime}) = \sum_ {i = 1} ^ {n} (t _ {i} - t _ {i - 1}) d \big (h (t _ {i - 1}), h ^ {\prime \prime} (t _ {i - 1}) \big) \\ \leqslant \sum_ {i = 1} ^ {n} (t _ {i} - t _ {i - 1}) \left(d \big (h (t _ {i - 1}), h ^ {\prime} (t _ {i - 1}) \big) + d \big (h (t _ {i - 1}), h ^ {\prime \prime} (t _ {i - 1}) \big)\right) \\ = d ^ {*} (h, h ^ {\prime}) + d ^ {*} (h, h ^ {\prime \prime}), \end{array}
\]

therefore \(d^{*}\) satisfies the triangle inequality.

This distance \(d^{*}\) is invariant under right translations (resp. left translations) if d is.

PROPOSITION 7. --- Suppose that G is a metrizable topological group. Then the topological group G* is metrizable. More precisely, if d is a bounded metric on G that defines its topology, then the topology of G* is defined by the metric d*.

Let \(d\) be a metric on G that defines its topology. Then, the map \(d'\) given by \(d'(h,h') = \inf (d(h,h'),1)\) is a bounded metric on G which also defines its topology (TG, IX, p. 3). It therefore suffices to prove that \(d^{*}\) defines the topology of \(G^{*}\) under the hypothesis that \(d\) is bounded.

Let V be a neighborhood of e in G, and let \(\varepsilon\) be a strictly positive real number. Let \(\delta\) be a strictly positive real number such that V contains the ball \(\mathrm{B}(e,\delta)\). Let \(h \in G^{*}\) and let \((t_{i})_{0 \leqslant i \leqslant n}\) be a sequence adapted to h.

Then,

\[
\begin{aligned} p_{\mathrm{V}}(h) & = \sum_{\substack{i = 1\\ h(t_{i - 1})\not\in\mathrm{V}}}^{n}(t_{i} - t_{i - 1})\\ & \leqslant \sum_{\substack{i = 1\\ d(h(t_{i - 1}),e)\geqslant \delta}}^{n}(t_{i} - t_{i - 1})\frac{d(h(t_{i - 1}),e)}{\delta}\\ & \leqslant \frac{d^{*}(h,e^{*})}{\delta}. \end{aligned}
\]

Therefore, the ball \(\mathrm{B}(e^{*},\varepsilon\delta)\) in \(G^{*}\) is contained in \(V_{\varepsilon}^{*}\).

Conversely, let \(\delta\) be a strictly positive real number and let \(\Delta\) be a strictly positive upper bound for the diameter of G. Let V be a neighborhood of \(e\) in G contained in the ball B \((e,\delta/2)\). For a function \(h\in\mathrm{G}^{*}\) and a sequence \((t_{i})_{0\leqslant i\leqslant n}\) adapted to \(h\), we have

\[
\begin{array}{l}d^{*}(h,e^{*}) = \sum_{i = 1}^{n}(t_{i} - t_{i - 1})d(h(t_{i - 1}),e)\\ = \sum_{\substack{i = 1\\ d(h(t_{i - 1}),e)\leqslant \delta /2}}^{n}(t_{i} - t_{i - 1})d(h(t_{i - 1}),e)\\ +\sum_{\substack{i = 1\\ d(h(t_{i - 1}),e) > \delta /2}}^{n}(t_{i} - t_{i - 1})d(h(t_{i - 1}),e)\\ \\ \leqslant \frac{\delta}{2} +\Delta \sum_{\substack{i = 1\\ h(t_{i - 1})\not\in\mathrm{V}}}^{n}(t_{i} - t_{i - 1})\\ \\ \leqslant \frac{\delta}{2} +\Delta p_{\mathrm{V}}(h). \end{array}
\]

The preceding inequality implies that, for every element \(h\) of \(\mathrm{V}_{\delta/2\Delta}^{*}\), we have \(d^{*}(h,e^{*})\leqslant\delta\). Consequently, every ball of \(\mathrm{G}^{*}\) for the distance \(d^{*}\) contains a neighborhood of the identity element.

Remarks. --- 2) Let d be a bounded metric defining the topology of G. When one equips the topological group \(G^{*}\) with the metric \(d^{*}\), the homomorphism \(\iota: G \to G^{*}\) is an isometry.

\begin{enumerate}
\def\labelenumi{\arabic{enumi})}
\setcounter{enumi}{2}
\tightlist
\item
  Suppose that G is a discrete topological group. The subgroup \(\iota(\mathrm{G})\) is then a discrete subgroup of \(G^{*}\). Indeed, the topology of G is defined by the metric d on G given by \(d(g,g') = 1\) if \(g \neq g'\) and \(d(g, g') = 0\) otherwise; the assertion then follows from the previous remark.
\end{enumerate}

THEOREM 2. --- Let G be a discrete topological group. Let G act on the right on G* by h · g = hι(g), and denote by B the quotient topological space G*/G.

The space G* is a covering of B, principal with group G; it is simply path-connected.

The space B is a metrizable topological space, path-connected, locally contractible, and its Poincaré group at every point is isomorphic to G.

Thus, the space B is a classifying space for the group G.

The group G* is contractible to its identity element (IV, p. 450, prop. 6). In particular, it is path-connected (III, p. 260) and simply path-connected (IV, p. 340).

The group \(\iota(\mathrm{G})\) is closed in \(\mathrm{G}^*\); therefore the space \(\mathrm{G}^*/\mathrm{G}\) is metrizable (TG, III, p. 13, prop. 18 and TG, IX, p. 25, prop. 4). It is also path-connected (III, p. 258, prop. 3). Since the group \(\iota(\mathrm{G})\) is discrete, it follows from Corollary 2 of I, p. 100, that \(\mathrm{G}^*\) is a principal G-bundle over B. The pointed covering \((\mathrm{G}^*, e^*)\) is then a universal covering of the pointed space B based at the image of \(e^*\); the Poincaré group of B (at each of its points) is isomorphic to G.

